% /root/monodromy, 2026-11-19 08:15 UTC. Complete proof independently audited by root and RD.
\subsection{Irreducible product Hasse divisors and all full images}
\label{subsec:product-hasse-tacnode}
\status{Proved. The complete Cartier, tacnode, small-degree jet and component-merging arguments were independently audited by the parent and the reduction-descent agent. This completes every degree for the fixed surface; it is not a classification for arbitrary ambient varieties.}

\begin{lemma}[A pointed Cartier calculation]
\label{lem:pointed-Cartier}
Let $C$ be a smooth projective curve in characteristic two with
one-dimensional Cartier kernel, generated by $\omega$.
If $P\in C$ satisfies $\omega(P)\ne0$, then Cartier induces a
surjection
\[
 H^0(C,\omega_C(2P))\longrightarrow H^0(C,\omega_C),
\]
and its kernel is $k\omega$. If $D$ is ordinary, the corresponding
map for $D$ is surjective with a one-dimensional kernel, whose
nonzero members have a nonzero double-pole coefficient at $P$.
\end{lemma}

\begin{proof}
There is no residue for a meromorphic differential with its only
possible pole at $P$, by the residue theorem. Cartier kills its
double-pole term, so the displayed target is correct.
The source has dimension $g(C)+1$ by Riemann--Roch.
For an ordinary curve, Cartier is already bijective on its
$g(C)$-dimensional holomorphic subspace, which proves the second
assertion.

For the first assertion, Serre duality and finite Frobenius
duality identify the displayed Cartier map with Frobenius
\[
 H^1(C,\mathcal O_C(-P))\longrightarrow
 H^1(C,\mathcal O_C(-2P)).
\]
The duality is the curve case of the local Cartier trace and
finite-duality formula in
\cite[Lemma~1.3.6 and Proposition~1.3.7, pp.~26--27]{BrionKumar2005};
its twists are exactly $F^*\mathcal O_C(-P)=\mathcal O_C(-2P)$.
The first space maps isomorphically to $H^1(C,\mathcal O_C)$,
while the kernel of the second space mapping to
$H^1(C,\mathcal O_C)$ is
$\mathfrak m_P/\mathfrak m_P^2$, the first-jet line.

Here is the boundary coefficient explicitly. Use an affine cover
with only one member $U_0$ containing $P$. A Cech representative
$c_{ij}$ of a nonzero Frobenius-kernel class satisfies
$c_{ij}^2=b_i-b_j$. The differentials $db_i$ glue, because
$d(c_{ij}^2)=0$, and give a nonzero global Cartier-killed
form $\omega$, up to scale. This also follows directly from
the local exact sequence
$0\to\mathcal O_C\xrightarrow{F}F_*\mathcal O_C
\xrightarrow{d} B^1_C\to0$: in a uniformizer, the exact
differentials are exactly those killed by Cartier.
The map on global constants is bijective, so the connecting
map identifies $H^0(B^1_C)$ with the Frobenius kernel on
$H^1(\mathcal O_C)$. In particular the glued differential is
nonzero.

Subtract the same constant from all the $b_i$ so that $b_0(P)=0$.
The Cech intersections avoid $P$, so the original cocycle also
represents its unique lift to $H^1(\mathcal O_C(-P))$.
Its Frobenius image is the boundary of
$b_0\bmod\mathfrak m_P^2\in\mathfrak m_P/\mathfrak m_P^2$.
The coefficient of a local parameter in this class is
$db_0(P)=\omega(P)$, up to the harmless sign of the connecting
homomorphism. Hence the lift of the one-dimensional Frobenius
kernel has nonzero image in this jet line. Any kernel of the
displayed twisted Frobenius would already lie in that original
kernel; it is therefore zero. Duality proves the surjectivity,
and the dimension calculation gives its Cartier kernel.
\end{proof}

\begin{lemma}[Smoothness at a mixed tacnodal union]
\label{lem:mixed-tacnode-Hasse}
Let $X=E\times E$ be the fixed ordinary surface above.
Let $C\in|L^a|$ and $D\in|L^e|$ be smooth curves without common
components, ordinary except that $C$ has a one-dimensional
Cartier kernel $k\omega$. Suppose they have a single simple
tangency at $P$, are otherwise transverse, and $\omega(P)\ne0$.
Then the full degree-$(a+e)$ Hasse divisor is smooth at $C\cup D$.
\end{lemma}

\begin{proof}
The normalization condition can be read directly from the
hypersurface dualizing generator, so no assertion about a smooth
normalization being a smooth family is required. In completed
coordinates the two branches have equations $y=0$ and
$y=\phi(x)$, with $\phi(x)=c x^2+O(x^3)$ and $c\ne0$.
The normalization ring is the pair of power-series rings, and
its subring for the union consists of pairs congruent modulo
$\phi(x)$, equivalently modulo $x^2$.
The adjunction generator $dx/(\partial(y(y-\phi))/\partial y)$
restricts on each branch to $dx/\phi(x)$, up to a sign, which
is immaterial in characteristic two. Multiplying by the
normalization subring shows that dualizing forms have at most
double poles, with matching double-pole coefficients and
matching residues. Conversely these two equalities are precisely
the congruence modulo $x^2$. At an ordinary node the same
calculation gives simple poles and opposite residues.

Cartier on dualizing forms acts componentwise on these rational
differentials. This follows either from the same local trace
formula or its compatibility with residues: the residue of a
Cartier image is the square root of the original residue.
The local matching conditions are preserved; the double-pole
term is killed and the simple-pole term becomes its square root.
A Cartier-killed dualizing form therefore has zero residues,
so it consists of Cartier-killed elements of
$H^0(\omega_C(2P))$ and $H^0(\omega_D(2P))$ with matching
double-pole coefficients. By
Lemma~\ref{lem:pointed-Cartier}, the first is a multiple
of the holomorphic $\omega$, while a nonzero element of the
second has nonzero polar coefficient. The matching condition
forces the second to be zero. Thus the Cartier kernel on
$C\cup D$ is exactly $k(\omega,0)$.

The exact adjunction sequence
$0\to\mathcal O_X\xrightarrow{s}L^{a+e}\to\omega_{C\cup D}\to0$
has an ordinary two-dimensional ambient Cartier cokernel on
global sections, because $H^1(X,\mathcal O_X)$ is ordinary.
Its kernel is the Cartier-unit line $ks$. Consequently its
Cartier kernels identify canonically: the one-dimensional
kernel just computed lifts uniquely, after adjusting by $ks$,
to
$\ker M(s)$, where $s=f_Cg_D$. Let $u$ be its generator.
Since its differential vanishes on $D$, the section is divisible
by $g_D$: $u=g_Da_0$ with $a_0\in H^0(L^a)$ and $a_0(P)\ne0$.

Choose local coordinates so $C$ is $y=0$ and $D$ has equation
$y+x^2+$ terms of higher order, allowing a nonzero coefficient
of $x^2$ and multiplying by a unit. The section $s$ has quadratic
term proportional to $y^2$, with no $xy$ term. The local Cartier
trace of a top differential at $P$ extracts its $xy$ coefficient.
Consequently the Frobenius-twisted evaluation at $P$ annihilates
$M(s)$ and spans its left kernel. Explicitly, for the matrix
convention with the output scalar square root removed, this is
the functional with coefficient vector $(e_i(P)^2)_i$ in a
local frame: $\mathcal C(sv)(P)^2=\operatorname{ev}^{[2]}_P M(s)v$.
This twist does not affect vanishing or the following nonzero
local coefficient. The section $u$ has linear term a nonzero multiple
of $y$. Choose a global deformation $\dot s$ with zero value at
$P$ and first jet $x$; very ampleness supplies it.
Then
\[
       \operatorname{ev}^{[2]}_P\bigl(M(\dot s)u\bigr)\ne0.
\]
The rank-one adjugate formula makes this a nonzero scalar
multiple of $dh_{a+e}(\dot s)$. Hence the Hasse divisor is
smooth at this union.
\end{proof}

\begin{lemma}[The degree-one jet exception is small]
\label{lem:product-small-degree-jets}
On the fixed product, $L^e$ for $e\ge2$ separates second jets at
one point and first jets at any two distinct points.
For $L$ itself, second jets are separated at a point
$P=(P_1,P_2)$ when neither $P_i$ is a flex in its plane cubic.
For every $Q\ne P$, the joint first-jet evaluation has rank at
least five; for fixed $P$, rank six fails only on the two fibres
through $P$ and a finite set of other points.
\end{lemma}

\begin{proof}
An elliptic line bundle of degree at least six interpolates
any subscheme of length at most four, since the residual line
bundle has positive degree and hence vanishing first cohomology.
Tensoring the interpolation maps on the two factors gives the
first assertion. For degree three, the map to length-three
jets at $P_i$ is an isomorphism unless
$\mathcal O_E(3P_i)\simeq\mathcal O_E(3O)$, precisely the flex
condition. Tensoring again gives second jets on the surface.

Write $V_i=H^0(E_i,\mathcal O(3O))^*$, let $e_P$ be the
one-dimensional evaluation direction in its projective plane,
and let $S_P\subset V_i$ be the two-dimensional vector tangent
plane to the embedded cubic. The dual of the surface first-jet
map at $P$ has image
\[
 A_P=e_{P_1}\otimes S_{P_2}+S_{P_1}\otimes e_{P_2},
 \qquad\dim A_P=3.
\]
Distinct points of a smooth plane cubic have distinct tangent
lines: a common tangent line would have total intersection
multiplicity at least four with a cubic.
If $P_1\ne Q_1$ and $P_2\ne Q_2$, then
\[
 A_P\cap A_Q\subseteq
 (S_{P_1}\cap S_{Q_1})\otimes(S_{P_2}\cap S_{Q_2}),
\]
whose dimension is one. A nonzero decomposable tensor in $A_P$
has either its first factor $e_{P_1}$ or its second factor
$e_{P_2}$; applying the same observation at $Q$ leaves only
\[
 e_{P_1}\in S_{Q_1},\ e_{Q_2}\in S_{P_2},
 \quad\text{or}\quad
 e_{Q_1}\in S_{P_1},\ e_{P_2}\in S_{Q_2}.
\]
For fixed $P$ these conditions permit only finitely many $Q$:
a tangent meets a cubic in a finite subscheme. For completeness,
tangents through a fixed point are finite for the specific
cubic here without any separability assumption on a Gauss map.
For $F=Y^2Z+XYZ+X^3+Z^3$, its partial derivatives are
\[
 F_X=YZ+X^2,\qquad F_Y=XZ,\qquad F_Z=Y^2+XY+Z^2.
\]
No nonzero constant linear combination is zero: compare the
coefficients of $X^2$, $Y^2$, and $XZ$. A nonzero quadratic
cannot vanish identically on the irreducible cubic. Hence for
every projective point $(a:b:c)$ the equation
$aF_X+bF_Y+cF_Z=0$ cuts a finite subscheme on $E$. Its points
are exactly those whose tangent contains $(a:b:c)$.
If, for example, $P_1=Q_1$ and $P_2\ne Q_2$, expanding in a
basis of $S_{P_1}$ beginning with $e_{P_1}$ shows
\[
 A_P\cap A_Q=e_{P_1}\otimes(S_{P_2}\cap S_{Q_2}),
\]
again of dimension one. This proves all stated rank bounds.
\end{proof}

\begin{theorem}[Uniform Hasse irreducibility]
\label{thm:product-Hasse-irred}
For the fixed surface and embedding, the degree-$d$ Hasse divisor
is geometrically irreducible for every $d\ge4$.
\end{theorem}

\begin{proof}
Work over an algebraic closure and put $e=d-3\ge1$.
Let $H_1,\ldots,H_t$ be the irreducible components of the
already reduced degree-three Hasse divisor in
$P_3=\mathbf P H^0(X,L^3)$.
Its general member on each $H_i$ is smooth, has stable loss one,
and has a one-dimensional Cartier kernel. Smoothness follows
because the singular-curve discriminant is irreducible and its
general member is ordinary: it contains an ordinary transverse
union of an ordinary degree-one and an ordinary degree-two
curve, so its nonempty ordinary open is dense.

For each $H_i$, choose a smooth general point $C$ and a point
$P\in C$ where its unique Cartier-killed differential is nonzero.
The tangent linear functional $dh_3|_C$ fails to be independent
of the first-jet functionals at at most one point of $C$.
Indeed first-jet evaluations at two distinct points are jointly
surjective by Lemma~\ref{lem:cubic-power-jets}; their dual spans
have zero intersection. Since $dh_3|_C\ne0$, it cannot belong
to both spans. We may therefore also impose this independence
at $P$, and require its two elliptic coordinates to avoid the
finite flex sets. These exclusions remove only finitely many
points of the smooth irreducible ample curve $C$.

Consider the incidence of $(C,P,D)$ over this open of pairs
$(C,P)$, with $D\in|L^e|$ tangent to $C$ at $P$.
The universal smooth curve over the irreducible open of $H_i$
is irreducible, since its fibres are connected smooth curves.
The two conditions on $D$ are independent linear first-jet
conditions, so the incidence is a projective bundle over an
irreducible base before removing open conditions.

The projection to $|L^e|$ is submersive on the chosen locus.
For precision, work on the affine cone over $H_i$: its tangent
space is $\ker dh_3$, and independence of $dh_3$ from the
first-jet functionals makes its map to the three-dimensional
first-jet space surjective. Projective scaling is harmless.
For an arbitrary infinitesimal change of $D$, move $P$ to
cancel its value change, using $dD(P)\ne0$. Then choose the
first jet of the change of $C$ inside $\ker dh_3$ to preserve
$C(P)=0$ and equality of tangent directions. These are exactly
the three incidence equations. The projection is consequently
dominant.

We also check nonemptiness of its good open, including $e=1$.
Fix a good $(C,P)$. If $e\ge2$,
Lemma~\ref{lem:product-small-degree-jets} gives the
usual independent conditions: second jets permit a nonzero
normal derivative and contact of exactly two at $P$;
a singularity of $D$ at another $Q$ imposes three further
conditions, while $Q$ varies in dimension two; another contact
with $C$ imposes two further conditions, while $Q\in C$ varies
in dimension one. Thus both bad loci are proper.

For $e=1$, second jets are still separated at the chosen $P$.
Outside the exceptional two fibres and finite set of
Lemma~\ref{lem:product-small-degree-jets}, the same
counts apply. On that exceptional locus, joint first-jet rank
is at least five. The two contact conditions at $P$ together
with three singularity conditions at $Q$ therefore have total
rank at least four; here $Q$ varies in dimension at most one.
Relative to the two fixed conditions, the bad locus again has
codimension at least one. For an additional contact at $Q\in C$,
the four conditions have rank at least three. The exceptional
$Q$ on $C$ are finite, since an irreducible ample curve contains
neither fibre through $P$. This bad locus is likewise proper.
These estimates prove that smooth $D$ with just the chosen
simple tangency and all other intersections transverse form a
nonempty open in the same irreducible incidence.

A nonempty open in an irreducible dominant incidence remains
dominant. It intersects the inverse image of the ordinary open
of $|L^e|$. Its image therefore contains a nonempty open of that
linear system. There are finitely many $H_i$, so a single general
smooth ordinary $D$ works for every $i$.

Fix this $D$, with equation $g$. On the multiplication subspace
$gP_3\subset\mathbf P H^0(L^d)$, the support of the Hasse divisor
is precisely
\[
                      H_1\cup\cdots\cup H_t\cup\Delta_D,
\]
as a divisorial statement. Here $\Delta_D$ is the irreducible
simple-contact divisor of
Theorem~\ref{thm:product-all-degree-transvections}.
This follows from the union stable-rank formula; exceptional
nonreduced members and members containing $D$ have codimension
at least two. The generic points of these displayed divisors
are distinct, and there is an ordinary member outside them.

For each $i$ the incidence just constructed gives a point of
$H_i\cap\Delta_D$ where the full Hasse divisor is smooth, by
Lemma~\ref{lem:mixed-tacnode-Hasse}. Any global Hasse component whose restriction contains either
$H_i$ or $\Delta_D$ contains this chosen intersection point,
because it contains that entire closed prime divisor on $gP_3$.
At a smooth point of the full Hasse divisor there is only one
such global component. Thus all components assigned to either
prime are the same. This argument allows arbitrary pullback
multiplicities and does not assume unique assignments in
advance.
Finally, every global Hasse component restricts nontrivially
to $gP_3$: a nonzero homogeneous polynomial of positive degree
cannot become a nonzero constant on this projective linear
subspace, and the subspace contains an ordinary point.
Thus its restriction contains one of the displayed prime
divisors. All have just been assigned to a single component,
so the full Hasse divisor is irreducible.
\end{proof}


\begin{corollary}[All full groups on the fixed ordinary product]
\label{thm:product-char2-all-full-groups}
Let $k$ be any perfect field of characteristic two, let
\[
 E:\ Y^2Z+XYZ+X^3+Z^3=0,\qquad
 X=E\times E,\qquad
 L=\mathcal O_E(3O)\boxtimes\mathcal O_E(3O),
\]
and use the fixed Segre embedding. For every $d\ge1$, the
maximal relative Hasse--Witt rank is $r_d=9d^2-1$. On the smooth
ordinary degree-$d$ hypersurface-section open, the entire
geometric and arithmetic images coincide and are
\[
 G_{{\rm geom},d}=G_{{\rm arith},d}=
 \begin{cases}
 \GL_{9d^2-1}(\mathbf F_2),&d\text{ odd},\\[2mm]
 \left\{\begin{pmatrix}A&B\\0&I_3\end{pmatrix}:
 A\in\GL_{9d^2-4}(\mathbf F_2),\quad
 B\in M_{(9d^2-4)\times3}(\mathbf F_2)\right\},&d\text{ even}.
 \end{cases}
\]
In even degree this is exactly the pointwise stabilizer of all
three independent invariant covectors, and its order is
$2^{3(9d^2-4)}|\GL_{9d^2-4}(\mathbf F_2)|$.
The representation is irreducible exactly in odd degree.
\end{corollary}

\begin{proof}
The generic rank and exact invariant covectors were proved in
Theorem~\ref{thm:product-char2-all-parities-fixed-quotient} and
the preceding product rank calculation. Degrees one and two
are Theorem~\ref{thm:product-char2-two-full-groups}; degree three
is Corollary~\ref{cor:product-degree-three-full-group}.
For $d\ge4$, Theorem~\ref{thm:product-Hasse-irred} discharges the
geometric hypothesis in
Theorem~\ref{thm:product-full-group-if-hasse-irreducible}.
Its proof supplies the entire linear group in odd degree and
all translations as well as the full linear kernel group in
even degree. The arithmetic and geometric equality persists
under any perfect extension of $\mathbf F_2$: the finite
geometrically connected covers giving the geometric images
remain connected after algebraically closed field extension,
and the three covectors in even degree are already defined
over $\mathbf F_2$. Finally the natural full linear module is
irreducible, while the even-degree common kernel of the three
covectors is a nonzero proper invariant subspace.
\end{proof}
